\documentclass[10pt]{article}
\providecommand{\FIGDIR}{fig}
\newcommand{\DIGESTDATE}{20 August 2026}
% ===================== PM Research Watch - shared preamble =====================
\usepackage[T1]{fontenc}
\usepackage[utf8]{inputenc}
\usepackage{amsmath,amssymb}
\usepackage{tgpagella}
\usepackage{tgheros}
\renewcommand{\sfdefault}{qhv}
\renewcommand{\ttdefault}{lmtt}
\usepackage{microtype}
\usepackage[english]{babel}

\usepackage[a4paper,top=17mm,bottom=17mm,left=17mm,right=17mm,
            headsep=6mm,footskip=10mm]{geometry}
\usepackage{graphicx}
\usepackage{xcolor}
\usepackage{colortbl}
\usepackage{tikz}
\usetikzlibrary{calc}
\usepackage[most]{tcolorbox}
\usepackage{titlesec}
\usepackage{fancyhdr}
\usepackage{booktabs}
\usepackage{array}
\usepackage{longtable}
\usepackage{enumitem}
\usepackage{multicol}
\usepackage{ragged2e}
\usepackage[hidelinks]{hyperref}

% ------------------------------------------------------------------ palette
\definecolor{Ink}    {HTML}{14262E}
\definecolor{Deep}   {HTML}{0F4C5C}
\definecolor{Teal}   {HTML}{1F7A8C}
\definecolor{Sky}    {HTML}{4FA3C4}
\definecolor{Sage}   {HTML}{6FA287}
\definecolor{Amber}  {HTML}{E3A008}
\definecolor{Clay}   {HTML}{B5643C}
\definecolor{Coral}  {HTML}{D1495B}
\definecolor{Violet} {HTML}{6A5B9E}
\definecolor{Slate}  {HTML}{4A6572}
\definecolor{Paper}  {HTML}{FBF9F5}
\definecolor{Mist}   {HTML}{EEEAE2}
\definecolor{Rule}   {HTML}{D9D3C9}

\pagecolor{Paper}
\color{Ink}

\hypersetup{colorlinks=true, linkcolor=Deep, urlcolor=Teal, citecolor=Teal,
            pdftitle={PM Research Watch}, pdfauthor={Automated daily digest}}

% ------------------------------------------------------------------ headings
\titleformat{\section}
  {\sffamily\bfseries\large\color{Deep}}{}{0pt}
  {}[{\vspace{1.2mm}\color{Amber}\titlerule[1.1pt]}]
\titlespacing*{\section}{0pt}{5.5mm}{2.8mm}

\titleformat{\subsection}
  {\sffamily\bfseries\normalsize\color{Teal}}{}{0pt}{}
\titlespacing*{\subsection}{0pt}{3.6mm}{1.4mm}

\setlist[itemize]{leftmargin=4.2mm, itemsep=1.1mm, topsep=1.3mm, parsep=0pt,
                  label=\textcolor{Teal}{\footnotesize$\blacktriangleright$}}

% ------------------------------------------------------------------ header/footer
\pagestyle{fancy}
\fancyhf{}
\renewcommand{\headrulewidth}{0pt}
\renewcommand{\footrulewidth}{0pt}
\fancyhead[L]{\sffamily\scriptsize\color{Slate}PM RESEARCH WATCH\;\textcolor{Rule}{\textbar}\;\DIGESTDATE}
\fancyhead[R]{\sffamily\scriptsize\color{Slate}Particulate matter monitoring \& health effects}
\fancyfoot[C]{\sffamily\scriptsize\color{Slate}\thepage}

% ------------------------------------------------------------------ components
% metric tile
\newcommand{\metric}[3]{%
  \begin{tcolorbox}[enhanced, colback=white, colframe=Rule, boxrule=0.5pt,
    arc=1.6mm, left=2.2mm, right=2.2mm, top=1.8mm, bottom=1.8mm,
    borderline west={1.7mm}{0pt}{#3}, nobeforeafter, valign=center,
    height=17mm, width=\linewidth]
    \RaggedRight
    {\sffamily\bfseries\fontsize{17}{17}\selectfont\color{#3}#1}\\[0.5mm]
    {\sffamily\fontsize{6.5}{7.4}\selectfont\color{Slate}#2}
  \end{tcolorbox}}

% pull-quote / key insight box
\newtcolorbox{keybox}[1][]{enhanced, colback=Mist, colframe=Mist,
  arc=1.6mm, left=3.2mm, right=3.2mm, top=2.4mm, bottom=2.4mm,
  borderline west={1.7mm}{0pt}{Amber}, #1}

% subtopic band
\newcommand{\band}[2]{%
  \vspace{3.2mm}\noindent
  \begin{tcolorbox}[enhanced, colback=#2, colframe=#2, arc=1.2mm,
    left=2.6mm, right=2.6mm, top=1.1mm, bottom=1.1mm, boxrule=0pt,
    width=\linewidth, nobeforeafter]
    {\sffamily\bfseries\small\color{white}#1}
  \end{tcolorbox}\vspace{1.6mm}}

% one paper entry: {short cite}{title}{journal}{doi}{body}
\newcommand{\paperentry}[5]{%
  \par\vspace{1.9mm}\noindent
  {\sffamily\bfseries\footnotesize\color{Ink}#1}\;%
  {\sffamily\fontsize{7}{8}\selectfont\color{Slate}#3}\\[0.4mm]
  {\footnotesize\itshape\color{Deep}#2}\\[0.6mm]
  {\footnotesize #4}\\[0.3mm]
  {\fontsize{6.6}{7.6}\selectfont\color{Slate}\texttt{doi:}\,\href{https://doi.org/#5}{\texttt{\detokenize{#5}}}}
  \par}

% caption style
\newcommand{\figcap}[1]{\par\vspace{1.0mm}
  {\fontsize{7.2}{8.6}\selectfont\color{Slate}\textsf{\textbf{Fig.}}\; #1\par}\vspace{1.5mm}}

\providecommand{\permil}{\textperthousand}

\setlength{\parindent}{0pt}
\setlength{\parskip}{0pt}
\renewcommand{\arraystretch}{1.18}


\begin{document}

% ==================================================================== MASTHEAD
\thispagestyle{empty}
\begin{tikzpicture}[remember picture, overlay]
  \fill[Deep] (current page.north west) rectangle
      ($(current page.north east)+(0,-67mm)$);
  \fill[Amber] ($(current page.north west)+(0,-67mm)$) rectangle
      ($(current page.north east)+(0,-68.4mm)$);
  \foreach \x/\y/\r in {22/12/0.9, 41/26/1.5, 63/9/0.7, 88/31/1.1, 112/17/1.9,
                        138/29/0.8, 159/13/1.3, 178/54/1.0, 196/21/1.6,
                        30/44/0.6, 74/57/1.2, 125/48/0.9, 168/6/0.8, 191/59/0.7,
                        52/18/1.0, 100/60/0.6, 148/20/1.1, 66/49/1.4, 205/45/0.8}
    \fill[Sky, opacity=0.30] ($(current page.north west)+(\x mm,-\y mm)$) circle (\r mm);
\end{tikzpicture}

\vspace*{4mm}
{\sffamily\fontsize{7.6}{9}\selectfont\color{Sky}%
 AUTOMATED DAILY LITERATURE DIGEST\;\textbullet\;PhD RESEARCH WATCH\par}
\vspace{2.2mm}
{\sffamily\bfseries\fontsize{27}{29}\selectfont\color{white}%
 Particulate Matter\par}
{\sffamily\fontsize{27}{29}\selectfont\color{Sky}%
 Monitoring \& Health Effects\par}
\vspace{5.0mm}
{\sffamily\fontsize{8.4}{10}\selectfont\color{white}%
 Issue 2026-08-20\;\;\textcolor{Amber}{\textbar}\;\;Entry window 20 August 2026\;\;%
 \textcolor{Amber}{\textbar}\;\;18 records in scope, 49 screened out\par}
\vspace{18mm}

% -------------------------------------------------------------- metric strip
\noindent
\begin{minipage}[t]{0.190\linewidth}\metric{18}{RECORDS IN SCOPE\\(49 SCREENED OUT)}{Deep}\end{minipage}\hfill
\begin{minipage}[t]{0.190\linewidth}\metric{14}{PARTICLE-SIDE\\RECORDS (78\%)}{Teal}\end{minipage}\hfill
\begin{minipage}[t]{0.190\linewidth}\metric{1}{RECORD WITH A\\HEALTH ENDPOINT}{Sky}\end{minipage}\hfill
\begin{minipage}[t]{0.190\linewidth}\metric{0}{EFFECT ESTIMATES\\EXTRACTABLE}{Amber}\end{minipage}\hfill
\begin{minipage}[t]{0.190\linewidth}\metric{4}{METADATA-ONLY\\DEPOSITS (22\%)}{Clay}\end{minipage}

\vspace{6mm}

% -------------------------------------------------------------- signal of day
\section*{\sffamily\bfseries\color{Deep}Signal of the day}
\vspace{-1.5mm}{\color{Amber}\rule{\linewidth}{1.1pt}}\vspace{2.5mm}

\begin{keybox}
{\sffamily\bfseries\small\color{Deep}A three-year, ten-node PurpleAir network anchored to
collocated reference monitors is now the pollution baseline for a capital city ---
the distributed low-cost architecture doing the job a regulatory network was never
built to do.}\\[1.4mm]
\footnotesize
Shlipak et al.\ (\textit{ACS EST Air}) report April 2022--March 2025 of continuous PM$_{2.5}$
from \textbf{10 PurpleAir sensors across Addis Ababa}, calibrated against Beta Attenuation
Monitors collocated at two contrasting U.S.\ Embassy sites, plus black carbon from
AethLabs MA350 microaethalometers at those same two sites. Network-mean PM$_{2.5}$ was
\textbf{30\,$\mu$g/m$^3$}, exceeding the WHO 24-hour guideline of 15\,$\mu$g/m$^3$ on
\textbf{99\% of days}; BC was 3.7 and 7.7\,$\mu$g/m$^3$ at the two BC sites, an
\emph{intra-urban} factor of two that a single-site regulatory monitor would have erased.
Aethalometer-model apportionment put \textbf{94\% of BC on fossil fuel} and 6\% on biomass
burning --- inverting the assumption that a sub-Saharan African capital's carbon is
predominantly biomass. Two design choices are worth stealing: the two-site BAM anchor
(rather than one) constrains the calibration across contrasting source regimes, and
tying the deployment to the forthcoming MAIA satellite gives the surface network a
defined validation role rather than leaving it as a standalone dataset. For distributed
low-cost PM work this is the strongest existing-practice reference published this window.
\end{keybox}

\vspace{3mm}

% -------------------------------------------------------------- exec brief
\section*{\sffamily\bfseries\color{Deep}What changed today}
\vspace{-1.5mm}{\color{Amber}\rule{\linewidth}{1.1pt}}\vspace{2.5mm}

\footnotesize
\begin{itemize}
\item \textbf{This is an atmospheric-science issue, not an epidemiology one.} 14 of 18
records measure particles rather than people, and \textbf{exactly one} carries a human
health endpoint. \textbf{No record in the window published a confidence-interval-bearing
risk estimate}, so the forest plot is omitted rather than populated with chamber
fold-changes and model fit statistics --- see \textit{Corpus analytics}.

\item \textbf{Two independent papers attack the same measurement gap: source resolution
without a mass-only monitor.} Shlipak et al.\ separate fossil from biomass BC by
aethalometer model in Addis Ababa; Saarikoski et al.\ (\textit{ACP}) show at a Helsinki
traffic site that relying on hydrocarbon-like OA alone to estimate traffic primary OA
\textbf{underestimates it by $\sim$50\%}, because a second, more oxygenated
traffic-related factor (distinct C$_2$H$_4$O$_2^+$, C$_2$H$_5$O$_2^+$, C$_3$H$_5$O$_2^+$
fractions) also peaks with the morning rush. Any source-apportionment product built on
an HOA-only convention inherits that bias.

\item \textbf{Wildfire smoke is being reframed as a joint-pollutant problem.}
Yu et al.\ (\textit{ES\&T}) fit a spatiotemporal Bayesian neural field over 2004--2023
CONUS at 5\,km to estimate daily probabilities of \emph{co-occurring} PM$_{2.5}$ and ozone
extremes, and find the amplified co-occurrence risk propagating into the Midwest and
Northeast on long-range transport. Two monitors, not one, is the implied network
requirement.

\item \textbf{Personal exposure and neighbourhood proxies diverge sharply, and the
divergence is quantified.} Zundel et al.\ (medRxiv preprint) report personal PM$_{2.5}$
of 28.2$\pm$20.0 vs 5.1$\pm$4.7\,$\mu$g/m$^3$ on smoke vs non-smoke days against
neighbourhood PurpleAir estimates of 87.6$\pm$80.2 vs 12.3$\pm$7.0 --- personal
concentrations were \emph{lower} than the outdoor proxy but rose fivefold, with a
within-person slope of only \textbf{2.3\,$\mu$g/m$^3$ per 10\,$\mu$g/m$^3$} of
neighbourhood PM$_{2.5}$. Four participants; treat the slope as an order of magnitude,
not an estimate.

\item \textbf{Aerosol phase state is becoming a tractable field measurement.}
Do et al.\ (\textit{ES\&T}) measured RH-dependent viscosity of ambient PM$_{2.5}$ by
poke-and-flow (6.2$\times$10$^4$--1.3$\times$10$^7$\,Pa\,s at 20--40\% RH) and propagated
it to nighttime N$_2$O$_5$ uptake coefficients of 0.01--0.07 --- up to an order of magnitude
below the liquid-particle assumption of 0.1 that chemical transport models still use.

\item \textbf{Metadata-only deposits are now a fixed 20--25\% tax.} 4 of 18 today
(Elsevier \textit{Atmos Environ} $\times$2, \textit{Atmos Res}, Springer Nature
\textit{JESEE}). Semantic Scholar and Europe PMC returned nothing for all four; they are
too new. They ship with an explicit \textsc{metadata only} marker and no inferred
quantity.

\item \textbf{Screening-side note.} 27 of the 49 exclusions came from a single journal
(MDPI \textit{Sensors}) whose ISSN was added to the Crossref sweep this run and returned
\textbf{zero} in-scope records --- it indexes sensing hardware of every kind, not airborne
particle instrumentation. Recorded for the next sweep-list revision.
\end{itemize}

\clearpage

% ==================================================================== ANALYTICS
\section*{\sffamily\bfseries\color{Deep}Corpus analytics}
\vspace{-1.5mm}{\color{Amber}\rule{\linewidth}{1.1pt}}\vspace{3mm}

\noindent\includegraphics[width=\linewidth]{\FIGDIR/f1_subtopics.png}
\figcap{Subtopic assignment of the 18 in-scope records. Assignment is single-label by
dominant contribution. \textit{Exposure assessment \& modelling} takes 10 of 18 and
\textit{Sensing, forecasting \& instrumentation} a further 4; the six clinically oriented
clusters between them hold three records. Four of the ten subtopic constants
(mechanistic toxicology, cardiovascular/metabolic, neuro/mental health, reproductive/%
developmental) are empty --- an unusually lopsided day, and the direct consequence of the
harvest's shape: PubMed's two axes returned 8 unique PMIDs of which 5 shipped, while the
Crossref-by-ISSN leg contributed the other 13.}

\vspace{2mm}
\noindent\includegraphics[width=\linewidth]{\FIGDIR/f2_design_endpoint.png}
\figcap{Left --- study architecture. Measurement campaigns lead at 6, followed by
metadata-only deposits at 4 and modelling/inventory work at 3; there is no cohort, no
case-crossover and no trial in this window. Right --- health endpoint by organ system.
Thirteen records explicitly measure particles rather than people and four more are
metadata-only deposits whose endpoint could not be extracted, leaving \textbf{exactly one
record with a stated human health endpoint} --- Zhang S et al., respiratory, a panel study
whose abstract reports direction and significance but no interval estimates.}

\vspace{2mm}
\noindent\includegraphics[width=\linewidth]{\FIGDIR/f3_metric_geography.png}
\figcap{Left --- particle metric. PM$_{2.5}$ alone is the exposure in 9 records; a further
5 concern unspeciated PM or emissions, and one each covers PM$_{10}$ alone, joint
PM$_{2.5}$/PM$_{10}$, optical properties (the Riyadh dust AOD/{\AA}ngstr\"om analysis) and
explicit composition/speciation. Right --- geographic provenance. Coverage is wider than
usual: North America, China and Europe hold 4 each, East Asia excluding China 3, and
there is one record each from sub-Saharan Africa (Addis Ababa) and the Middle East
(Riyadh). Both of the latter are monitoring-network papers, not health studies.}

\vspace{2mm}
\noindent\includegraphics[width=\linewidth]{\FIGDIR/f6_lifecourse.png}
\figcap{Life-course windows addressed, counted over the four records that involve human
participants at all. Working-age adults account for three (the Detroit wearable panel,
the Guangzhou healthy-adult panel, the Oregon resident survey); one paediatric record
(Seoul atopic dermatitis and asthma, metadata-only) and one adolescent participant
inside the Detroit panel complete the count. \textbf{No record this window addressed the
preconception or in-utero window, and none targeted adults over 65} --- the two stages
that ordinarily dominate this figure.}

\vspace{2mm}
\noindent\includegraphics[width=\linewidth]{\FIGDIR/f4_heatmap.png}
\figcap{Subtopic $\times$ study-architecture cross-tabulation. The mass sits in a single
block --- measurement campaigns and modelling within \textit{Exposure assessment} --- and the
health-side rows are almost empty. Counted against the store --- a sensing- or
exposure-side record carrying a non-null human health endpoint --- that
instrumentation/epidemiology bridge appears in \textbf{11 of the 19 preceding August
issues} and is \textbf{absent today}. On an 18-record day dominated by one leg of the
harvest, that absence is a sampling artefact of the window, not evidence about the field.
(The 19 August issue put this count at 12; recounting under an explicit null-endpoint
rule gives 11.)}

\vspace{2mm}
\begin{keybox}
{\sffamily\bfseries\footnotesize\color{Deep}No forest plot this issue}\\[1.2mm]
{\fontsize{7.2}{8.8}\selectfont
Every quantity published in this entry window is an absolute concentration or deposition
rate (Addis Ababa PM$_{2.5}$ and BC; Swiss microplastic deposition, 881
MPs\,m$^{-2}$\,d$^{-1}$, 95\% CI 562--1199 at the urban site), a chamber or model
fold-change (0.75--1.3$\times$ organic-aerosol mass enhancement; $\sim$40-fold NO$_x$
enhancement under Cl$\bullet$ cycling), or a fit statistic ($R^2=0.955$, RMSE 9.05 for the
Linfen AQI model). \textbf{None of these is a risk estimate}, and pooling them onto the
shared logarithmic risk axis would corrupt the weekly rollup, which aggregates this
field across issues. \texttt{EFFECTS} is therefore empty for 2026-08-20 by decision, not
by omission, and figure~5 is not rendered.\par}
\end{keybox}

\clearpage

% ==================================================================== DIGEST
\section*{\sffamily\bfseries\color{Deep}Paper-level digest}
\vspace{-1.5mm}{\color{Amber}\rule{\linewidth}{1.1pt}}\vspace{1mm}

{\fontsize{7.4}{9}\selectfont\color{Slate}Ordered by cluster. Each entry gives the
methodological core and the finding that matters, not an abstract paraphrase.\par}

\band{Sensing, forecasting \& instrumentation\hfill 4 records}{Amber}

\paperentry{Shlipak et al.\ 2026}
{Spatial and temporal trends of ambient PM$_{2.5}$ and source-apportioned black carbon in Addis Ababa}
{ACS EST Air}
{Three continuous years (Apr 2022--Mar 2025) from the MAIA surface network: 10 PurpleAir
units across the city calibrated against collocated Beta Attenuation Monitors at two
contrasting U.S.\ Embassy sites (Central, Jacros), plus AethLabs MA350 microaethalometers
at those two sites. Network mean PM$_{2.5}$ \textbf{30\,$\mu$g/m$^3$}; BC 3.7 and
7.7\,$\mu$g/m$^3$; the WHO 24\,h guideline of 15\,$\mu$g/m$^3$ was exceeded on
\textbf{99\% of days}. Diurnal peaks at morning traffic and under nighttime inversion,
strongest in the wet season; highest concentrations west of the city centre.
Aethalometer-model apportionment attributes \textbf{94\% of BC to fossil combustion},
6\% to biomass --- the opposite of the usual prior for the region. Directional and
back-trajectory analysis raises long-range shipping transport from the Gulf of Aden as a
regional contributor. Limitation: only two sites carry a reference anchor and BC, so the
calibration transfer across the other eight nodes is assumed rather than demonstrated.
The single most directly transferable deployment design in this window.}
{10.1021/acsestair.6c00140}

\paperentry{Zundel et al.\ 2026 (preprint)}
{Wearable personal PM$_{2.5}$ monitoring during wildfire smoke episodes in Metro Detroit}
{medRxiv}
{AirBeam3 wearables on four participants (one adolescent, three adults) during southeast
Michigan smoke episodes, with neighbourhood outdoor PM$_{2.5}$ taken as the mean of the
three nearest PurpleAir units and smoke days identified from state advisories.
\textbf{Personal 28.2$\pm$20.0 vs 5.1$\pm$4.7\,$\mu$g/m$^3$} on smoke vs non-smoke days;
neighbourhood \textbf{87.6$\pm$80.2 vs 12.3$\pm$7.0}. Within participants, each
10\,$\mu$g/m$^3$ of neighbourhood PM$_{2.5}$ mapped to only \textbf{2.3\,$\mu$g/m$^3$}
personal --- personal exposure is both attenuated and poorly tracked by the outdoor proxy,
yet still put $>$25\% of participant-days over the EPA 24\,h standard. $n=4$ and no
formal uncertainty on the slope: this is a feasibility demonstration, and the attenuation
factor should not be reused as a calibration constant. Preprint, not peer reviewed.}
{10.64898/2026.07.30.26359253}

\paperentry{Clark \& Bashar 2026}
{How IoT enables the Protective Action Decision Model during a wildfire event}
{Risk Anal}
{Survey of 1{,}200 Oregon residents (Jun--Aug 2021) about their responses to the 2020
fire season, framed against the Protective Action Decision Model and the availability of
hyper-local real-time PM$_{2.5}$ from IoT devices. The finding that matters for network
designers is negative: residents were \textbf{overwhelmed by the data and reported
inadequate guidance} on how to act on it. Retrospective self-report a year after the
event, and the survey does not establish who actually had sensor access --- so this
constrains interface and messaging design rather than measuring sensor efficacy. Still
the rare paper in this watch that treats the last mile of a dense network as the
problem.}
{10.1111/risa.70319}

\paperentry{Wei et al.\ 2026}
{Stacked Transformer-XGBoost model with SHAP interpretability for AQI prediction in Linfen}
{iScience}
{Transformer temporal encoder stacked with XGBoost through temporally consistent
out-of-fold stacking, on 2013--2024 air quality and meteorological series from a heavily
industrialised Chinese city. Reported \textbf{$R^2=0.955$, RMSE 9.05, MAE 2.91}; SHAP
attributes the variance to PM$_{2.5}$, NO$_2$, lagged AQI and seasonal terms. The
performance claim is against conventional ML and DL baselines on the same series, with
no external site held out, so it is an in-sample-city result. The out-of-fold stacking
discipline is the part worth borrowing --- naive stacking on autocorrelated air quality
series leaks badly.}
{10.1016/j.isci.2026.116396}

\band{Exposure assessment \& modelling\hfill 10 records}{Amber}

\paperentry{Yu et al.\ 2026}
{Probabilistic assessment of PM$_{2.5}$ and ozone co-occurring extremes under wildfire smoke}
{Environ Sci Technol}
{Spatiotemporal Bayesian neural field (ST-BayesNF) integrating satellite smoke-coverage
products, ground PM$_{2.5}$ and ozone, and meteorological reanalysis to give daily
probabilities of \emph{joint} extremes at 5\,km over the contiguous US, 2004--2023. The
model discriminates co-extreme from non-co-extreme days and is most sensitive to smoke
density interacting with boundary-layer height and near-surface temperature; predicted
co-occurrence risk rises during smoke events not only in the western US but increasingly
in the Midwest and Northeast under long-range transport. No health outcome is modelled ---
the health premise (co-occurrence raises risk beyond either pollutant alone) is imported
from prior literature and not tested here.}
{10.1021/acs.est.6c04403}

\paperentry{Saarikoski et al.\ 2026}
{Role of hydrocarbons and oxygenated organic species in traffic-derived aerosol}
{Atmos Chem Phys}
{Four Aerodyne AMS campaigns (2018--2024) at a Helsinki traffic site with PMF source
apportionment. Beyond hydrocarbon-like OA peaking at the morning rush, a second
traffic-related factor (TrOA) peaks in the morning but stays elevated longer, with mass
spectra between HOA and biomass-burning OA and distinct C$_2$H$_4$O$_2^+$ ($m/z$ 60),
C$_2$H$_5$O$_2^+$ (61) and C$_3$H$_5$O$_2^+$ (73) fractions. \textbf{Using HOA alone to
estimate traffic primary OA underestimates it by $\sim$50\%.} The authors are explicit
that TrOA's origin --- atmospheric processing versus direct emission from modern
vehicles --- is unresolved, which is exactly the ambiguity that makes the 50\% figure a
bound rather than a correction factor. Directly relevant to any speciated or
source-resolved sensing product.}
{10.5194/acp-26-11817-2026}

\paperentry{Do et al.\ 2026}
{Field-constrained diurnal viscosity of urban PM$_{2.5}$ and implications for N$_2$O$_5$ uptake}
{Environ Sci Technol}
{Poke-and-flow measurement with fluid-dynamics inversion on PM$_{2.5}$ collected in Ansan,
South Korea (summer 2024): \textbf{6.2$\times$10$^4$--1.3$\times$10$^7$\,Pa\,s at
20--40\% RH}, exceeding 10$^8$\,Pa\,s below $\sim$10\% RH. Pooled with Seoul and Beijing
data into a Vogel--Tammann--Fulcher temperature--RH parametrisation and pushed through an
ML surrogate over hourly meteorology for 14 megacities. Predicted viscosity is
systematically lower at night everywhere (RH plasticisation beating nocturnal cooling),
giving nighttime N$_2$O$_5$ uptake coefficients of \textbf{0.01--0.07} against the
conventional liquid assumption of 0.1. Weakness stated by the authors: the
parametrisation rests on 24\,h composition, so no time-resolved chemistry enters the
diurnal prediction it is being used to make.}
{10.1021/acs.est.6c04811}

\paperentry{Wang Q et al.\ 2026}
{Quantitative analysis of aerosol acidity and its driving factors in the Guanzhong area}
{Atmos Chem Phys}
{One year of water-soluble inorganic ions in PM$_{2.5}$ on the semi-arid Guanzhong Plain,
with thermodynamic pH reconstruction and sensitivity-based driver attribution. Annual
mean pH \textbf{3.8$\pm$1.0} (winter $>$ spring $>$ summer $>$ autumn), shifting from 2--5
on clean days toward 3--6 on polluted days. Temperature contributes 22.3--33.8\% of pH
variation across all seasons, NH$_x$ 11.4--44.9\% (dominant in autumn/winter), sulfate
8.5--10.8\%; Ca$^{2+}$ is a spring-specific driver. RH shows a non-monotonic effect with
an inflection at 60--85\% RH, attributed to deliquescence-driven liquid water content.
Attribution rests on a thermodynamic model rather than measured pH, which is the standing
caveat for this whole literature.}
{10.5194/acp-26-11785-2026}

\paperentry{Kommula et al.\ 2026}
{Comparable organic aerosol mass enhancement during photochemical and dark oxidation of biomass burning emissions}
{Environ Sci Technol}
{Chamber aging of emissions from savannah grass, savannah wood and boreal forest surface
fuel under varied combustion conditions, comparing daytime photochemical with nighttime
dark oxidation. OA is 82--99\% of PM$_1$ in the fresh emissions. Both aging routes yield
\textbf{comparable net OA mass}, which PMF traces to more efficient primary-OA loss during
daytime offsetting more efficient gas-phase oxidation. The resulting mass enhancement of
\textbf{0.75--1.3$\times$} matches field observations and closes a long-standing
lab--field discrepancy. Applies to low-NO$_x$ remote conditions; the authors say so, and
that is where the proposed model simplification is and is not safe.}
{10.1021/acs.est.5c15841}

\paperentry{Albugami 2026}
{Dust event characteristics, transport pathways and source regions over Riyadh}
{Atmosphere}
{34 dust events over 2024--2025 identified from AERONET by a dual optical criterion
(AOD$_{500}\geq$0.50 and 440--870\,nm {\AA}ngstr\"om exponent $\leq$0.50), cross-checked
against NOAA ISD present-weather and visibility and traced with HYSPLIT back-trajectories,
PSCF and CWT. \textbf{82\% fell in March--May}; mean {\AA}ngstr\"om exponent 0.28, mean
fine-mode fraction 0.26. Surface present-weather codes confirmed 26 of 34; the other
8 (24\%) were classified as elevated dust invisible at the surface station --- a clean
demonstration that column and surface instruments disagree by design, not by error.
Source--receptor analysis points to the Mesopotamian Basin, with the northern
Rub' al-Khali secondary. Single-author, single-station, two years.}
{10.3390/atmos17080802}

\paperentry{Ashta et al.\ 2026}
{Wet and dry atmospheric deposition of microplastics across a Swiss site gradient}
{Atmos Chem Phys}
{Four-weekly wet and dry deposition sampling over one year (May 2024--May 2025) at urban
(Zurich), suburban (D\"ubendorf), rural (Magadino, Payerne) and mountainous (Chaumont)
sites, with focal-plane-array $\mu$-FTIR over the 20--215\,$\mu$m range and an explicit
measurement-uncertainty treatment. Urban deposition \textbf{881 MPs\,m$^{-2}$\,d$^{-1}$
(95\% CI 562--1199)} and \textbf{53\,$\mu$g\,m$^{-2}$\,d$^{-1}$ (17--107)}; the other four
sites are similar to each other at 249--331 MPs\,m$^{-2}$\,d$^{-1}$ (140--478). Scaled to
land use, 219\,t or 3.8$\times$10$^{14}$ particles per year below 2000\,m nationally,
excluding tyre wear. The 20\,$\mu$m lower cut is the binding limitation: the
inhalable fraction sits mostly below it.}
{10.5194/acp-26-11803-2026}

\paperentry{Wang T et al.\ 2026}
{Chlorine radical cycling on photoactive particulate matter amplifies photolytic NO$_x$ release}
{Environ Sci Technol}
{Bench photochemistry showing chloride on photoactive PM surfaces oxidised by
photogenerated holes to Cl$\bullet$, which accelerates conversion of nitrite intermediates
to gaseous NO$_x$ while chloride is regenerated. \textbf{$\sim$40-fold NO$_x$ enhancement}
under the model conditions, robust across humidity, chloride salt, substrate (including
soot) and nitrogen source (ammonium, 4-nitrophenol). Rate scales with photoactive-matter
loading and PM composition. Extrapolated to \textbf{0.7--6.2\% of NO$_2$ variability} in
the study region --- a small number derived from a large laboratory enhancement, and the
gap between the two is the result's main uncertainty.}
{10.1021/acs.est.6c06768}

\paperentry{Naghizade et al.\ 2026}
{Multi-class organic marker profiling of urban PM$_{10}$ in Strasbourg}
{Atmos Environ}
{\textsc{metadata only} --- seasonal distribution of exhaust, non-exhaust and
biomass-burning tracers in northeastern France. No abstract in the Crossref deposit;
Semantic Scholar and Europe PMC both returned nothing on 20 August. No quantity is
inferred here. Flagged for the next issue that can reach the full record.}
{10.1016/j.atmosenv.2026.122306}

\paperentry{Chuang et al.\ 2026}
{Simulating PM$_{2.5}$ in southern Taiwan under high-pressure return flow (Kao-Ping Experiment)}
{Atmos Res}
{\textsc{metadata only} --- a KPEx modelling study of a synoptic regime that traps
PM$_{2.5}$ over the Kao-Ping region. No abstract in the deposit and no indexed full text
on 20 August. No quantity inferred.}
{10.1016/j.atmosres.2026.109286}

\clearpage

\band{Respiratory \& allergic\hfill 1 record}{Sage}

\paperentry{Zhang S et al.\ 2026}
{Respiratory and cardiovascular effects of ozone modified by PM$_{2.5}$}
{Environ Sci Technol}
{Two complementary designs testing one hypothesis: that 30-day mean PM$_{2.5}$
\emph{modifies} the acute cardiorespiratory response to ozone, and that unmodelled
modification is part of why the ozone literature is inconsistent. In a healthy-adult
panel, short-term ozone was associated with significantly larger decrements in FEV$_1$,
FEF$_{25-75}$ and PEF, and larger rises in systolic and diastolic blood pressure, under
\textbf{PM$_{2.5}$ above 22\,$\mu$g/m$^3$}. A meta-regression across published studies
agreed: study areas with higher annual PM$_{2.5}$ showed larger ozone effects on FEV$_1$
and blood pressure. \textbf{The abstract reports direction and significance but no
interval estimates}, which is why this record contributes nothing to the forest plot.
Panel size is not stated in the abstract either. Mechanistically the claim is priming
rather than co-exposure confounding, and the meta-regression cannot separate the two ---
areas with high PM$_{2.5}$ differ in many ways. Still the day's one health record, and
the modification framing is directly relevant to designing multi-pollutant sensor nodes.}
{10.1021/acs.est.6c05604}

\band{Occupational \& indoor exposure\hfill 1 record}{Clay}

\paperentry{Zhao et al.\ 2026}
{CFD-based prediction of air change rate and PM in a dynamically occupied operating room}
{Indoor Air}
{Class 10{,}000 operating room in Dalian (Chinese GB 50333-2013), combining field
measurement with CFD across 45 orthogonal scenarios, Sobol sensitivity, and three
optimisation routes (Bayesian, multi-objective, response surface) to calibrate a
non-uniformity coefficient and fit ACH--occupancy--PM. Response surface gave the lowest
prediction error (\textbf{RMSE 0.032, MAE 0.024}); multi-objective optimisation was the
most stable across runs ($C_v$ 0.005). The premise is sound --- fixed as-built ACH
standards ignore occupancy dynamics --- but the model is calibrated on one room in one
climate, and the PM metric is cleanroom particle count rather than a health-relevant
mass fraction.}
{10.1155/ina/5838379}

\band{Burden, policy \& mitigation\hfill 1 record}{Coral}

\paperentry{Anastasopolos et al.\ 2026}
{Mortality attributed to ambient air pollution by anthropogenic emission sector in Canada}
{Atmos Environ}
{\textsc{metadata only} --- a comparative health-burden analysis run at multiple
geographic scales with sectoral attribution. Nine authors, Environment and Climate Change
Canada style. No abstract in the Crossref deposit and nothing indexed elsewhere on
20 August; the sector-resolved burden numbers, which are the point of the paper, are
therefore not reportable yet. Flagged for a later issue.}
{10.1016/j.atmosenv.2026.122310}

\band{Other clinical endpoints\hfill 1 record}{Slate}

\paperentry{Lee et al.\ 2026}
{Short-term ambient VOC mixtures, paediatric atopic dermatitis and asthma in Seoul}
{J Expo Sci Environ Epidemiol}
{\textsc{metadata only} --- short-term effects and attributable burden of ambient volatile
organic compound \emph{mixtures} on two paediatric outcomes. Included at the edge of scope:
the exposure is VOCs rather than particles, but mixture-burden methods on a shared urban
monitoring backbone are directly transferable, and the same Seoul network carries
PM$_{2.5}$. No abstract available on 20 August, so neither the mixture method nor the
burden fraction is reportable.}
{10.1038/s41370-026-00965-5}

\vspace{4mm}

% ==================================================================== TABLE
\section*{\sffamily\bfseries\color{Deep}Corpus register}
\vspace{-1.5mm}{\color{Amber}\rule{\linewidth}{1.1pt}}\vspace{2.5mm}

{\fontsize{6.7}{8.4}\selectfont
\setlength{\tabcolsep}{2.4pt}
\begin{longtable}{@{}>{\RaggedRight}p{25mm}>{\RaggedRight}p{41mm}>{\RaggedRight}p{27mm}%
>{\RaggedRight}p{23mm}>{\RaggedRight}p{18mm}>{\RaggedRight}p{15mm}@{}}
\toprule
\textsf{\textbf{First author}} & \textsf{\textbf{Focus}} & \textsf{\textbf{Journal}} &
\textsf{\textbf{Design}} & \textsf{\textbf{Metric}} & \textsf{\textbf{Region}}\\
\midrule
\endfirsthead
\toprule
\textsf{\textbf{First author}} & \textsf{\textbf{Focus}} & \textsf{\textbf{Journal}} &
\textsf{\textbf{Design}} & \textsf{\textbf{Metric}} & \textsf{\textbf{Region}}\\
\midrule
\endhead
\multicolumn{6}{@{}l@{}}{\cellcolor{Mist}\textsf{\textbf{\textcolor{Deep}{Neuro / mental health}}}}\\[0.6mm]
\href{https://doi.org/10.1001/jamanetworkopen.2026.25015}{Dales} & Wildfire-sourced PM$_{2.5}$ and ED visits for migraine and other primary headache & \textit{JAMA Netw Open} & Case-crossover & PM$_{2.5}$ (wildfire) & Canada\\
\href{https://doi.org/10.1093/ageing/afag215}{Glover} & Umbrella meta-review of urban environmental exposures and cognitive health & \textit{Age Ageing} & Umbrella review / meta-analysis & PM$_{2.5}$ & Global\\
\href{https://doi.org/10.1016/j.nsa.2026.107019}{Goossens} & In utero air pollution exposure and pre-/postnatal brain development & \textit{Neurosci Appl} & Narrative review & PM$_{2.5}$ / PM$_{10}$ / PM$_{\mathrm{coarse}}$ & Global\\
\href{https://doi.org/10.1016/j.isci.2026.116617}{Ma} & PM$_{2.5}$-induced neurotoxicity and depression: exposure risks and mechanisms & \textit{iScience} & Narrative review & PM$_{2.5}$ + components & Global\\
\href{https://doi.org/10.3760/cma.j.cn112150-20250804-00752}{Pan} & Maternal PM exposure and offspring neurodevelopment: effects and mechanisms & \textit{Zhonghua Yu Fang Yi Xue Za Zhi} & Narrative review & PM (bulk) & China\\
\addlinespace[0.9mm]
\multicolumn{6}{@{}l@{}}{\cellcolor{Mist}\textsf{\textbf{\textcolor{Teal}{Cardiovascular \& metabolic}}}}\\[0.6mm]
\href{https://doi.org/10.3760/cma.j.cn112338-20251120-00843}{Huang} & PM$_{2.5}$ and dyslipidaemia in rural Chinese children; dietary effect modification & \textit{Zhonghua Liu Xing Bing Xue Za Zhi} & Cross-sectional (multistage) & PM$_{2.5}$ & China\\
\href{https://doi.org/10.32394/pe/218966}{Kazmierska} & Air pollution and cardiovascular disease: review of current data & \textit{Przegl Epidemiol} & Narrative review & PM$_{2.5}$ / PM$_{10}$ & Global\\
\href{https://doi.org/10.1007/s00125-026-06814-2}{Xi} & Air pollution, built environment and temperature vs. prediabetes phenotypes & \textit{Diabetologia} & Prospective cohort & PM$_{2.5}$ & Germany\\
\href{https://doi.org/10.1161/JAHA.125.046917}{Yang} & Short-term PM$_{2.5}$ and hospitalisation/ED use in atrial fibrillation & \textit{J Am Heart Assoc} & Case-crossover & PM$_{2.5}$ (wildfire-influenced) & USA\\
\addlinespace[0.9mm]
\multicolumn{6}{@{}l@{}}{\cellcolor{Mist}\textsf{\textbf{\textcolor{Sky}{Reproductive \& developmental}}}}\\[0.6mm]
\href{https://doi.org/10.1016/j.jes.2025.10.043}{Chen} & Prenatal air pollution, fetal growth trajectory phenotypes and child neurodevelopment & \textit{J Environ Sci (China)} & Prospective cohort & PM$_{2.5}$ / PM$_{10}$ & China\\
\href{https://doi.org/10.1016/j.ecoenv.2026.120566}{Huang} & Joint prenatal bisphenol + air-pollutant mixture exposure and preschool IQ & \textit{Ecotoxicol Environ Saf} & Prospective cohort (mixtures) & PM$_{2.5}$ / PM$_{10}$ & China\\
\href{https://doi.org/10.1007/s40572-026-00552-8}{Peterson} & Ambient air pollution, PAHs and pubertal development: critical review & \textit{Curr Environ Health Rep} & Critical review & PM$_{2.5}$ components & Global\\
\href{https://doi.org/10.1016/j.envres.2026.125313}{Xie} & Long-term PM$_{2.5}$ and timing of pubertal development in boys and girls & \textit{Environ Res} & Prospective cohort & PM$_{2.5}$ & China\\
\addlinespace[0.9mm]
\multicolumn{6}{@{}l@{}}{\cellcolor{Mist}\textsf{\textbf{\textcolor{Sage}{Respiratory \& allergic}}}}\\[0.6mm]
\href{https://doi.org/10.1007/s00484-026-03289-z}{Bayar Muluk} & Climate change and otolaryngology: airway disease impacts and emission reduction & \textit{Int J Biometeorol} & Narrative review & PM$_{2.5}$ & Global\\
\href{https://doi.org/10.4103/mgr.MEDGASRES-D-25-00164}{Liao} & Short-term air pollution, respiratory symptoms and lung function in asthma & \textit{Med Gas Res} & Panel study & PM$_{2.5}$ / PM$_{10}$ & China\\
\href{https://doi.org/10.3390/epidemiologia7040097}{Sanchez de Prada} & Ambient PM and asthma-coded admissions in children 0-3 y, Spain & \textit{Epidemiologia} & Ecological panel & PM$_{2.5}$ / PM$_{10}$ & Spain\\
\href{https://doi.org/10.1016/j.jhazmat.2026.143076}{Yang} & Allergenic vs. toxic risk of airborne seed hairs compared with ambient PM$_{2.5}$ & \textit{J Hazard Mater} & Field measurement & PM$_{2.5}$ + bioaerosol & China\\
\addlinespace[0.9mm]
\multicolumn{6}{@{}l@{}}{\cellcolor{Mist}\textsf{\textbf{\textcolor{Violet}{Mechanistic toxicology}}}}\\[0.6mm]
\href{https://doi.org/10.1016/j.jes.2025.12.074}{Ding} & PM$_{2.5}$ promotes vascular calcification via miR-27b-5p/HES1 signalling & \textit{J Environ Sci (China)} & In vitro + organoid + murine & PM$_{2.5}$ & China\\
\href{https://doi.org/10.1007/s00360-026-01698-5}{Maraschi} & Osmoregulatory disruption and metal bioaccumulation from settleable atmospheric PM & \textit{J Comp Physiol B} & Controlled exposure (ecotox) & Settleable APM & Brazil\\
\href{https://doi.org/10.1016/j.jes.2025.11.018}{Wu} & Macrophage NLRP3 inflammasome in real-ambient PM-induced hepatic glucose dysregulation & \textit{J Environ Sci (China)} & Real-ambient murine + co-culture & Real-ambient PM & China\\
\addlinespace[0.9mm]
\multicolumn{6}{@{}l@{}}{\cellcolor{Mist}\textsf{\textbf{\textcolor{Clay}{Occupational \& indoor}}}}\\[0.6mm]
\href{https://doi.org/10.3389/fmedt.2026.1796828}{Chen} & Nested suction instrument to cut intraoperative surgical-smoke PM$_{2.5}$ & \textit{Front Med Technol} & Bench / device experiment & PM$_{2.5}$ (surgical smoke) & China\\
\href{https://doi.org/10.1093/annweh/wxag055}{Cole} & Determinants of respirable crystalline silica exposure in Australian tunnelling & \textit{Ann Work Expo Health} & Observational exposure study & Respirable dust / RCS & Australia\\
\href{https://doi.org/10.1016/j.jhazmat.2026.143081}{Ding} & Stage-resolved hair Zn-isotope dynamics as a firefighting combustion-exposure tracer & \textit{J Hazard Mater} & Repeated-measures biomarker & Combustion-derived PM & China\\
\href{https://doi.org/10.3390/jox16040119}{Ramsden} & Xenobiotic hazards in aircraft cabin air & \textit{J Xenobiot} & Review + risk economics & Ultrafine particles & Global\\
\addlinespace[0.9mm]
\multicolumn{6}{@{}l@{}}{\cellcolor{Mist}\textsf{\textbf{\textcolor{Amber}{Sensing, forecasting \& instrumentation}}}}\\[0.6mm]
\href{https://doi.org/10.1016/j.jenvman.2026.130515}{Shi} & Adaptive spatiotemporal graph transformer for multi-site multi-step PM$_{2.5}$ forecasting & \textit{J Environ Manage} & Deep-learning model & PM$_{2.5}$ & China / India\\
\href{https://doi.org/10.3390/bios16070380}{Tsuchiya} & Seasonal and indoor-outdoor characteristics of fluorescent aerosol particles in offices & \textit{Biosensors} & Real-time sensor campaign & Bioaerosol / coarse & Japan\\
\addlinespace[0.9mm]
\multicolumn{6}{@{}l@{}}{\cellcolor{Mist}\textsf{\textbf{\textcolor{Amber}{Exposure assessment \& modelling}}}}\\[0.6mm]
\href{https://doi.org/10.1029/2025GH001621}{Solorzano-Araque} & Multiscale exposure modelling: LOTOS-EUROS + mobility surveys + particle toxicity & \textit{GeoHealth} & CTM + mobility integration & PM$_{2.5}$ / PM$_{10}$ & Colombia\\
\addlinespace[0.9mm]
\multicolumn{6}{@{}l@{}}{\cellcolor{Mist}\textsf{\textbf{\textcolor{Coral}{Burden, policy \& mitigation}}}}\\[0.6mm]
\href{https://doi.org/10.1002/hsr2.72886}{Karamzad} & Ischaemic heart disease burden in MENA 1990-2021; PM attributable fraction 28.2\% & \textit{Health Sci Rep} & GBD secondary analysis & Ambient PM & MENA\\
\href{https://doi.org/10.1371/journal.pone.0347989}{Rees} & HEPA filtration to reduce respiratory infection in care homes & \textit{PLoS One} & RCT process evaluation & Filtered PM & UK\\
\href{https://doi.org/10.1016/j.jes.2026.01.009}{Shan} & Policy-driven emission reductions and post-pandemic rebound, high-resolution inventories & \textit{J Environ Sci (China)} & Emission inventory & PM emissions & China\\
\addlinespace[0.9mm]
\multicolumn{6}{@{}l@{}}{\cellcolor{Mist}\textsf{\textbf{\textcolor{Slate}{Other clinical endpoints}}}}\\[0.6mm]
\href{https://doi.org/10.3390/dj14070440}{Lupita} & PM$_{2.5}$, preventive-care capacity and caries/edentulism burden across the EU & \textit{Dent J} & Ecological panel & PM$_{2.5}$ & EU-27\\
\href{https://doi.org/10.1016/j.jdin.2026.06.008}{Park} & Long-term PM exposure and risk of androgenetic alopecia & \textit{JAAD Int} & Nationwide cohort & PM$_{2.5}$ / PM$_{10}$ & South Korea\\
\href{https://doi.org/10.3171/2026.3.PEDS2615}{Sun} & Environmental exposure and area deprivation vs. paediatric hydrocephalus outcomes & \textit{J Neurosurg Pediatr} & Registry cohort + geospatial & PM$_{2.5}$ (satellite-derived) & USA\\
\addlinespace[0.9mm]
\bottomrule
\end{longtable}}


\vspace{2mm}
\begin{keybox}
{\sffamily\bfseries\footnotesize\color{Deep}Provenance and method}\\[1.2mm]
{\fontsize{7.2}{8.8}\selectfont
Entry window \textbf{20 August 2026}, contiguous with the 19 August issue.
\textbf{PubMed} (E-utilities and connector) was queried on \texttt{[Date - Entry]} along a
health axis and a monitoring/sensing axis as separate queries: 6 and 3 local hits, 2 and 5
connector hits, \textbf{8 unique PMIDs}, of which 5 shipped. \textbf{Europe PMC} returned
0 --- and the topic-free control query (\texttt{CREATION\_DATE:[2026-08-20 TO 2026-08-20]},
no topic terms) also returned 0 records of any kind, so this is index lag, not a broken
leg. \textbf{OpenAlex} returned HTTP 429 for the 18th consecutive run and remains dead.
\textbf{arXiv} returned nothing inside the 3-day recency screen. \textbf{Crossref by
ISSN} carried the issue: 5 deposits across the 8 standing journals plus 54 across 18
supplementary ISSNs added this run, of which 13 shipped.
Subtopic, design, particle-metric and geography labels are assigned from title, abstract,
keywords and MeSH terms; assignment is single-label by dominant contribution and is a
judgement, not a controlled vocabulary. Four records shipped as \textsc{metadata only}:
the publisher deposited no abstract and neither Semantic Scholar nor Europe PMC had
indexed them; no quantity is inferred for those records.
\textbf{No effect estimate with a confidence interval was published in this window}, so
the forest plot is omitted --- see the note in \textit{Corpus analytics}. 49 records were
screened out with a logged reason, 27 of them from one supplementary ISSN (MDPI
\textit{Sensors}) that returned no in-scope work. All DOI links resolve to the publisher
of record and were checked by \texttt{check\_dois.py} before this issue was built.
Content derived from PubMed and Crossref metadata; abstracts remain the copyright of
their respective publishers.\par}
\end{keybox}

\end{document}
